% SGA 6, Exposición VIII: fragmento español alineado con la fuente.
% Alcance exacto: sección 4 completa, fr.tex, líneas 11649--11812.
% SHA-256 de la fuente francesa completa:
% 0D2158B9FE3B466309A92165727A0B191422E05FBAA0DACD5806E8322C54BA1A.
% Termina inmediatamente antes de la siguiente unidad estructural.
% Las superficies matemáticas se conservan diplomáticamente; las anomalías
% detectadas se notifican por separado para su resolución canónica.

\section*{4. Teorema de Riemann--Roch: caso de una inmersión cerrada regular}

\textbf{4.1.} Sea \(i:Y\to X\) una inmersión cerrada regular. El teorema de
Riemann--Roch para \(i\) es esencialmente la relación VII 4.4, que permite
calcular las clases de Chern de \(i_*(x)\) a partir de los \(\gamma^p(N,x)\).
Más precisamente, si se introduce el \(\lambda\)-anillo
\(K^{\bullet}(Y)_{\mathbb Q N}\) definido en V 5.5, el homomorfismo de
\(K^{\bullet}(Y)_{\mathbb Q N}\) en \(K^{\bullet}(X)_{\mathbb Q}\) deducido de
\(i_*\) es un \(\lambda\)-homomorfismo por VII 2.8 y VII 4.3. El teorema de
Riemann--Roch se reduce entonces a una fórmula análoga relativa al homomorfismo
(aditivo) \(K^{\bullet}(Y)\longrightarrow K^{\bullet}(Y)_{\mathbb Q N}\)
definido por \(y\rightsquigarrow(0,y)\), es decir, a una relación de naturaleza
formal, válida para todo \(\lambda\)-anillo. No desarrollaremos aquí esta
demostración, para evitar todavía más formalismo sobre los \(\lambda\)-anillos
(cf. \(\RRR\) I 1.5), y retomaremos el método de [1], que utilizaba los esquemas
obtenidos por explosión (método ya empleado en la demostración de VII 4.3 y que
es también el de la parte final de \(\RRR\)).

Como se indicó en 3.9, se puede restringir al caso de una inmersión regular de
codimensión constante \(d\) que satisface la condición de VII 4.1; el caso de
una inmersión regular general se deduce después de 3.8. Retomaremos en este
parágrafo las notaciones de VII 3.1. Se tiene, pues, el cuadrado cartesiano
\[
\begin{tikzcd}[column sep=large,row sep=large]
Y' \arrow[r,"j"] \arrow[d,"g"'] & X' \arrow[d,"f"]\\
Y \arrow[r,"i"'] & X ,
\end{tikzcd}
\]
donde \(X'\) es el \(X\)-esquema obtenido efectuando la explosión de \(Y\). Se
supone además que \(X\) es cuasicompacto y posee un haz amplio.

\textbf{Proposición 4.2.} El morfismo \(f\) es proyectivo, de intersección
completa y de dimensión relativa virtual nula. Se tiene además la relación
\[
 f_{\mathrm{gr}}\circ f^{\mathrm{gr}}=\operatorname{Id}_{\Gr K^{\bullet}(X)_{\mathbb Q}} .
\]
La primera afirmación se deduce inmediatamente de (VII 1.8). La segunda
relación es consecuencia de la relación
\[
 f_*\circ f^*=\operatorname{Id}_{K^{\bullet}(X)},
\]
demostrada en (VII 3.6).

\textbf{Proposición 4.3.} Para todo \(y\in K^{\bullet}(Y)_{\mathbb Q}\), se
tiene la relación
\[
 f^{\mathrm{gr}}\circ i_{\mathrm{gr}}(y)
 =
 j_{\mathrm{gr}}\bigl(g^{\mathrm{gr}}(y)\cdot c^{d-1}(\check F)\bigr),
\]
donde \(c^{d-1}\) denota la clase de Chern \((d-1)\)-ésima.

Esta relación está contenida en VII 4.8, pues no es sino
\[
 v^{\mathrm{gr}}\circ u^{\mathrm{gr}}=0 .
\]

\textbf{Proposición 4.4.} Sea \(X\) un esquema cuasicompacto. Si \(E\) es un
haz localmente libre de tipo finito sobre \(X\), de clase \(E\) en
\(K^{\bullet}(X)\) y de rango \(n\), se tiene
\[
 \ch(\lambda_{-1}(E))=c^n(\check E)(\Todd(\check E))^{-1}.
\]

Los dos miembros de la relación que se debe demostrar son multiplicativos: es
claro para \(\ch\circ\lambda_{-1}\) y para \(\Todd\); también es cierto para
\(c^n\), que es la clase de
\((-1)^{\varepsilon(E)}\lambda_{-1}(E)\). Por consiguiente, pasando al fibrado
de banderas de \(E\), se puede suponer que \(E\) es una suma de clases de haces
invertibles y reducirse a demostrar la fórmula para un haz invertible. Sea
\(\xi\) la clase de tal haz. Se tiene entonces
\[
 \ch(\lambda_{-1}(\xi))=\ch(1-\xi)=1-\ch(\xi).
\]
Como \(\widetilde c(\xi)=(1,1+c^1(\xi))\), se obtiene, por la definición de
\(\ch\),
\[
 \ch(\lambda_{-1}(\xi))=1-\exp(c^1(\xi)).
\]
Por otra parte, \(\xi\) tiene aumentación \(1\); la clase de Chern que
interviene en la fórmula 4.4 relativa a \(\xi\) es, pues, \(c^1\). Teniendo en
cuenta la definición de \(\Todd\) y 3.5, la relación que se debe demostrar se
escribe
\[
 1-\exp(c^1(\xi))
 =
 -c^1(\xi)\cdot\left(\frac{-c^1(\xi)}{1-\exp(c^1(\xi))}\right)^{-1},
\]
y queda así verificada.

\textbf{Observación 4.4.1.} También es posible deducir 4.4 de un enunciado
puramente algebraico correspondiente; cf. \(\RRR\) I 1.3.

\textbf{4.5.} Demostraremos ahora que, para demostrar el teorema de
Riemann--Roch en el caso considerado en este parágrafo, nos podemos reducir al
caso en que \(Y\) es un divisor de \(X\) y en que el elemento al cual se aplica
la fórmula es la imagen inversa de un elemento de \(K^{\bullet}(X)\).

Por 3.2, la fórmula que se debe demostrar,
\[
 \ch(i_*(x))=
 i_{\mathrm{gr}}\bigl(\ch(x)\cdot\Todd(\check T_i)\bigr)
\]
es equivalente a la fórmula
\[
 f^{\mathrm{gr}}\bigl(\ch(i_*(x))\bigr)
 =
 f^{\mathrm{gr}}\circ i_{\mathrm{gr}}\bigl(\ch(x)\cdot\Todd(\check T_i)\bigr).
\]
Teniendo en cuenta la funtorialidad de \(\ch\) y 4.3, también puede escribirse
\[
 \ch\bigl(f^*(i_*(x))\bigr)
 =
 j_{\mathrm{gr}}\bigl(g^{\mathrm{gr}}\bigl(\ch(x)\cdot\Todd(\check T_i)\bigr)\cdot c^{d-1}(\check F)\bigr),
\]
es decir, por VII 3.4,
\[
 \ch\bigl(j_*(g^*(x)\cdot\lambda_{-1}(F))\bigr)
 =
 j_{\mathrm{gr}}\bigl(\ch(g^*(y))\cdot\Todd(g^*(\check T_i))\cdot c^{d-1}(\check F)\bigr).
\]

Ahora bien, por definición se tiene
\[
 T_i=-\cl(J/J^2)=-N,\qquad
 T_j=-\cl(J'/J'^2)=-L.
\]
Por otra parte, la sucesión exacta
\[
0\longrightarrow F\longrightarrow g^*(N)\longrightarrow \mathcal O_{Y'}(1)\longrightarrow0
\]
da la relación
\[
 g^*(T_i)=-(F+L),
\]
de donde se deduce
\[
 \Todd(g^*(\check T_i))=\Todd(\check T_j)\cdot(\Todd(\check F))^{-1}.
\]
Teniendo en cuenta 4.4, la relación que se debe demostrar se escribe, pues,
\[
 \ch\bigl(j_*(g^*(x)\cdot\lambda_{-1}(F))\bigr)
 =
 j_{\mathrm{gr}}\bigl(\ch(g^*(x)\cdot\lambda_{-1}(F))\cdot\Todd(\check T_j)\bigr).
\]
Poniendo \(z=g^*(x)\cdot\lambda_{-1}(F)\), esta relación no es sino el teorema
de Riemann--Roch aplicado a la inmersión regular de codimensión \(1\)~\(j\), y
al elemento \(z\) de \(K^{\bullet}(Y')\).

Además, como se ha supuesto que \(i\) satisface la hipótesis de (VII 4.1),
\(\lambda_{-1}(F)\) es divisible por \((1-L)\). Se puede, pues, encontrar
\(z'\) tal que \(z=z'(1-L)\), o también \(z=j^*(j_*(z'))\) (Exp. VII 2.7). Se
reduce así al caso en que \(Y\) es un divisor y a un elemento \(x\) tal
que existe \(y\), con \(x=i^*(y)\).

\textbf{4.6.} Hay que demostrar, por tanto, la relación
\[
 \ch\bigl(i_*(i^*(y))\bigr)
 =
 i_{\mathrm{gr}}\bigl(\ch(i^*(y))\cdot\Todd(\check T_i)\bigr).
\]
Ahora bien, la fórmula de proyección da
\[
 i_*(i^*(y))=y\cdot i_*(1),
\]
y de la sucesión exacta
\(0\longrightarrow J\longrightarrow\mathcal O_X\longrightarrow\mathcal O_Y\longrightarrow0\),
donde \(J\) es aquí invertible, se deduce
\[
 i_*(1)=1-J,
\]
de donde
\[
 \ch(i_*(i^*(y)))=\ch(y)\cdot\ch(1-J).
\]
Como \(J\) es invertible, se tiene
\(\widetilde c(J)=(1,1+c^1(J))\), y por tanto
\(\ch(J)=\exp(c^1(J))\).
Por consiguiente, se tiene
\[
 \ch(i_*(i^*(y)))=\ch(y)\cdot(1-\exp(c^1(J))).
\]

Por otra parte, la funtorialidad de \(\ch\) y la fórmula de proyección 3.7
implican
\[
 i_{\mathrm{gr}}\bigl(\ch(i^*(y))\cdot\Todd(\check T_i)\bigr)
 =
 \ch(y)\cdot i_{\mathrm{gr}}(\Todd(\check T_i)).
\]
Ahora bien, \(T_i=-\cl(J/J^2)=-i^*(J)\). Por tanto,
\[
 i_{\mathrm{gr}}(\Todd(\check T_i))
 =
 i_{\mathrm{gr}}\bigl(i^{\mathrm{gr}}(\Todd(-\check J))\bigr)
 =
 i_{\mathrm{gr}}(1)\cdot\Todd(-\check J).
\]
Como \(i_*(1)=1-J\), \(i^{\mathrm{gr}}_*(1)\) es la imagen en
\(\Gr^1K^{\bullet}(X)\) de \(1-J\), y por tanto es \(-c(J)\). Por otra parte,
se tiene
\[
 \Todd(-\check J)=(\Todd(\check J))^{-1}
 =
 \frac{1-\exp(c^1(J))}{-c^1(J)}.
\]
Se obtiene finalmente
\[
 i_{\mathrm{gr}}\bigl(\ch(i^*(y))\cdot\Todd(\check T_i)\bigr)
 =
 \ch(y)\cdot(1-\exp(c^1(J))),
\]
lo que demuestra la relación buscada.
